% Copyright (c) 2026 Geovana Neves. Licensed under CC BY 4.0 (https://creativecommons.org/licenses/by/4.0/). See LICENSE-AND-AUTHORSHIP.md.
%
\section{Rotor coefficients}

\begin{frame}{Rotor coefficients describe one rotor}
\begin{itemize}
  \item \code{J}, \code{CT}, \code{CQ}, \code{CP}, \code{ETA}, \code{ETAW}:
        one table per rotor, every column suffixed with that rotor's alias.
  \item They are meaningful for one rotor. Summing several rotors mixes
        different diameters and speeds in the normalisation.
\end{itemize}
\[
  J = \frac{V}{|n|\,D}, \quad
  C_T = \frac{T}{\rho n^2 D^4}, \quad
  C_Q = \frac{Q}{\rho n^2 D^5}, \quad
  C_P = 2\pi\,C_Q\,\mathrm{sign}(n), \quad
  \eta = \frac{J\,C_T}{C_P}
\]
\begin{itemize}\small
  \item \(n\) is the signed speed in revolutions per \textbf{second}: its
        magnitude sets \(J\), its sign multiplies \(C_P\). These are the
        classical airscrew coefficients [6, 7].
\end{itemize}
\end{frame}

\begin{frame}{Tip and helical Mach numbers}
\[ \Omega = \frac{2\pi\,\mathrm{RPM}}{60}, \qquad
   M_\mathrm{tip} = \frac{\Omega R}{a}, \qquad
   M_\mathrm{hel} = \frac{\sqrt{V^2 + (\Omega R)^2}}{a} \]
\begin{itemize}\small
  \item Stated for each rotor an \code{unsteady\_rotor} or
        \code{qsteady\_rotor} row turns, for the rotor of a \code{steady} row
        that states \code{RPM}, and for each actuator disc a row names. \(R\) is half the rotor's diameter, or the
        disc's tip radius; \(V\) and \(a\) are the point's own resolved free
        stream and speed of sound. At \(V = 0\), \(M_\mathrm{hel} =
        M_\mathrm{tip}\).
  \item The rotor table ends with \code{MTIP\_<alias>} and
        \code{MHEL\_<alias>}; \code{plan.json} and the run record carry both
        under \code{rotor\_mach}. A point whose condition does not resolve
        reads \code{NA} in both and keeps its row.
  \item Both are geometric: \(V\) is the free stream alone, without the
        rotor's induced velocity, which near hover adds to the speed the tip
        sees. At \(M_\mathrm{hel} \ge 1\) the tip is sonic, beyond the
        linearised compressibility of a panel method, and the plan warns
        once, naming every such point, its rotor and its \(M_\mathrm{hel}\)
        [1, 10].
\end{itemize}
\end{frame}

\begin{frame}{An unsteady rotor row uses the shared history window}
\begin{itemize}
  \item A steady point's row comes from the loads export.
  \item An unsteady point's row averages the plots history over the
        matrix row's window, shared with the unsteady POLAR and reductions.
  \item The loads export of an unsteady run is the last time step,
        one instant of a cycle.
\end{itemize}
\end{frame}

\begin{frame}{The rotor history is a global group over its own families}
\begin{itemize}\small
  \item Use the rotor's six components \code{FX, FY, FZ, MX, MY, MZ}
        in global \code{MRP}, over exactly its general and blade families.
  \item Find them through a \code{[[plots.groups]]} entry with
        \code{frame = "MRP"} and those exact families, whatever its name.
  \item If the pproc plots none, the run adds \code{ROTOR\_<ALIAS>}.
  \item A group in the rotor's own frame cannot supply this history:
        its forces turn with the rotor and its moments are already about the hub.
  \item A group that only shares the alias's name but covers other families
        is not a source either.
\end{itemize}
\end{frame}

\begin{frame}{A windowed rotor row never substitutes an instant}
\begin{itemize}\small
  \item A point whose history does not cover the window, or lacks the required
        columns, is left out and named in \code{products.json}.
  \item The unsteady POLAR applies the same exclusion rule.
  \item The rotor table's manifest entry states \code{source} and \code{window}.
\end{itemize}
\end{frame}

\begin{frame}{\code{ETAW}: the full force vector in wind axes}
\begin{enumerate}
  \item Take the full rotor-frame force vector, $[F_x,F_y,F_z]_{\mathrm{rotor}}$.
  \item Carry it to the airframe body frame with the \textbf{transpose}
        of the rotor-to-body rotation.
  \item Carry it to wind axes with the \textbf{AIAA rotation}, using
        both alpha and beta.
  \item Take its X component, \code{Fx\_W}.
\end{enumerate}
\end{frame}

\begin{frame}{\code{ETAW}: a dimensionless efficiency}
\[
  \mathrm{CTW}=\frac{F_{x,W}}{\rho n^2D^4},
  \qquad
  \mathrm{ETAW}=\frac{J\,\mathrm{CTW}}{\mathrm{CP}}.
\]
\begin{itemize}\small
  \item Wind-axis force replaces thrust, normalised exactly as thrust is.
        \code{CTW} enters where \code{CT} enters the efficiency.
  \item \code{ETAW} reduces to \code{ETA} when the shaft lies along the stream.
  \item No wind-axis force supplied means \code{NA}.
  \item Two vector rotations preserve the off-axis force components.
        The cosine of a scalar angle does not.
\end{itemize}
\end{frame}

\begin{frame}{A static point has only a real zero advance ratio}
\begin{itemize}\small
  \item Every rotor coefficient is \code{NA} at a static point except
        \code{J}, which is \code{0.00000}.
  \item The export gives coefficients normalised by the run's dynamic pressure.
        At $V=0$, that pressure is zero: the real thrust has already been
        divided away before the package sees it.
  \item No rotor coefficient can be recovered from that static export.
  \item A hover figure of merit cannot be offered either: it needs a force
        the run does not state.
\end{itemize}
\end{frame}
